% Inputtable analytic fragment. Wrapper: moment-main.tex.
\section{Using the joint moment circle in initial acquisition}
\label{v6moment:section}

Let $P$ be centered with variance one and finite third absolute moment.
Write
\[
 \rho=\mathbb E|X|^3,\quad \gamma=\mathbb EX^3,\quad
 \beta=\mathbb E[X|X|],\quad m=\mathbb EX_+=\mathbb EX_->0,
 \quad \sigma=\sqrt{10}.
\]
All moments in this note belong to this one marginal. The moment-circle
inequality is already present in the manuscript's Lemma
\texttt{ms:physical-defect}; the new step here is its explicit inversion
and use in the initial scalar feedback. No new independent bound for the
Berry--Esseen constant is claimed.

\begin{lemma}[Same-law moment constraint]\label{v6moment:circle}
For $a_\pm=\mathbb EX_\pm^2=(1\pm\beta)/2$ and
$b_\pm=\mathbb EX_\pm^3=(\rho\pm\gamma)/2$,
\begin{equation}\label{v6moment:sign}
 mb_\pm\ge a_\pm^2,\qquad m^2\le a_+a_-,\qquad
 \gamma^2+4m^2\le\rho^2.
\end{equation}
An atom at zero is allowed.
\end{lemma}
\begin{proof}
Both signs have positive probability, since the law is centered and
nondegenerate. Cauchy--Schwarz applied on each sign to
$X_\pm^{1/2},X_\pm^{3/2}$ gives $a_\pm^2\le mb_\pm$.
If $p_\pm=P(\pm X>0)$, a second application gives
$m^2\le p_\pm a_\pm$. Thus
\[
 m^2(1/a_++1/a_-)\le p_++p_-\le1.
\]
As $a_++a_-=1$, this proves $m^2\le a_+a_-$, without assigning a zero
atom a false sign. Multiplying the first two bounds gives
\[
 b_+b_-\ge(a_+a_-)^2/m^2\ge m^2.
\]
Since $4b_+b_-=\rho^2-\gamma^2$, the circle follows. In particular
$2m\le1$ and $2m\le\sqrt{\rho^2-\gamma^2}$.
\end{proof}

For completeness, use the actual gap sum $J$ of the manuscript and its
conserved defect
\[
 D=(3+\sigma)\rho-\gamma-3J.
\]
Its gap bound $J\le\rho+2m$ remains valid before a central gap is
established. Indeed, on each gap lying on one side of zero,
$(b-a)^2\le2\int_a^b|x|\,dx$. Summing with its nonnegative gap weight
and using $2\int |x|\tau(x)\,dx=\rho$ gives the one-sided contribution.
If a central gap $(-A,B)$ exists, its additional contribution is $2mAB$.
Since $(X+A)(X-B)\ge0$ on support,
$0\le\mathbb E[(X+A)(X-B)]=1-AB$. Thus $AB\le1$ and the bound follows.
If zero belongs to support there is no central gap. Consequently
\begin{equation}\label{v6moment:defect}
 D\ge\sigma\rho-\gamma-6m
 \ge\sigma\rho-\gamma-3\sqrt{\rho^2-\gamma^2}.
\end{equation}
The first relation is an inequality, not an alternative definition of $D$.

\begin{lemma}[Inverting a defect ceiling]\label{v6moment:inverse}
Suppose $1\le l\le\rho\le u$, $D\le d$ with $d\ge0$, and
$\sigma_l\le\sqrt{10}$. Put $a_l=\sigma_l-d/l$. Define
\[
 Y(a)=
 \begin{cases}
 -1,&a\le-1,\\
 \displaystyle\frac{a-3\sqrt{10-a^2}}{10},&-1<a\le\sqrt{10}.
 \end{cases}
\]
Then $\gamma/\rho\ge Y(a_l)$. A rational lower bound $y_l\le Y(a_l)$
with $y_l\ge-1$ yields the valid cell bound
\begin{equation}\label{v6moment:cellgamma}
 \gamma\ge
 \begin{cases}l y_l,&y_l\ge0,\\u y_l,&y_l<0.\end{cases}
\end{equation}
\end{lemma}
\begin{proof}
Let $y=\gamma/\rho\in[-1,1]$. Dividing
\eqref{v6moment:defect} by $\rho$ gives
\[
 f(y):=y+3\sqrt{1-y^2}\ge\sigma-D/\rho
 \ge\sigma_l-d/l=a_l.
\]
The case $a_l\le-1$ imposes no new lower bound. On
$[-1,1/\sqrt{10}]$, $f$ increases continuously from $-1$ to
$\sqrt{10}$; its derivative in the interior is
$1-3y/\sqrt{1-y^2}$ and changes sign only at $1/\sqrt{10}$.
Solving $f(y)=a$ on this increasing branch gives
$10y^2-2ay+a^2-9=0$ and the stated smaller root.
It has $a-y\ge0$, including $-1\le a\le0$, because there
$\sqrt{10-a^2}\ge-3a$. Hence squaring introduces no spurious branch.
Every $y$ below this root has $f(y)<a$, proving the lower bound.
Finally multiplying $\gamma/\rho\ge y_l$ by the positive $\rho$
requires the endpoint $l$ for $y_l\ge0$ and $u$ for $y_l<0$.
\end{proof}

In exact computation the square root in $Y$ is rounded upward, so the
negative root term rounds $Y$ downward. The branch $a_l\le-1$ keeps
the whole interval. The implementation also intersects with the two
previously justified bounds $\gamma\ge-g$ and
$\gamma\ge\sigma_l l-3-d$. Thus
\begin{align}
 g_l&=\max\{-g,\sigma_l l-3-d,\text{bound in
 \eqref{v6moment:cellgamma}}\},\label{v6moment:gl}\\
 \rho-\gamma&\le
 \min\{2u,(1-\sigma_l)l+3+d,u-g_l\}.
 \label{v6moment:diff}
\end{align}
These inequalities retain the same actual $m$, $\rho$ and $\gamma$;
they do not independently optimize over moments from different laws.
